% Part: first-order-logic
% Chapter: axiomatic-deduction
% Section: provability-quantifiers

% verification of properties of provability needed for maximally
% consistent sets in the completeness chapter.

\documentclass[../../../include/open-logic-section]{subfiles}

\begin{document}

\olfileid{fol}{axd}{qpr}

\olsection{\usetoken{S}{derivability} ও পরিমাণসূচক}

\begin{explain}
  পূর্ণতা উপপাদ্যের জন্য আরও প্রয়োজন যে স্বতঃসিদ্ধমূলক নিষ্পাদন থেকে
  এই অংশে প্রতিষ্ঠিত~$\Proves$-সংক্রান্ত তথ্যগুলি পাওয়া যায়।
\end{explain}

\begin{thm}
\ollabel{thm:strong-generalization} $c$ যদি এমন কোনো !!a{constant} হয়
যা $\Gamma$ বা~$!A(x)$-এ আসে না এবং $\Gamma \Proves !A(c)$ হয়,
তবে $\Gamma \Proves \lforall[x][!A(x)]$।
\end{thm}

\begin{proof}
নিঃসরণ উপপাদ্য অনুসারে $\Gamma \Proves \ltrue \lif !A(c)$। যেহেতু
$c$, $\Gamma$ বা~$\ltrue$-এ আসে না, তাই $\Gamma \Proves \ltrue
\lif \lforall[x][!A(x)]$ পাই। সত্য-ধ্রুবকটি স্বতঃসিদ্ধ হওয়ায় মোডাস পোনেন্স থেকে
$\Gamma \Proves \lforall[x][!A(x)]$ পাই।
% BN-SRC-065 TARGET-CORRECTION-BEGIN
\par\smallskip\noindent\textnormal{\small\textbf{উৎস-সংশোধন (BN-SRC-065):} হিমায়িত ইংরেজি প্রমাণটি প্রথম ও তৃতীয় সূত্রে কনফিগারযোগ্য সত্য-ধ্রুবক $\ltrue$ ব্যবহার করলেও নতুন-ধ্রুবকের শর্তে কাঁচা $\top$ লেখে। বাংলা পাঠে একই সূত্রকে সর্বত্র $\ltrue$ দিয়ে চিহ্নিত করা হয়েছে।} % BN-SRC-065 TARGET-CORRECTION-END
% BN-SRC-071 TARGET-CORRECTION-BEGIN
\par\smallskip\noindent\textnormal{\small\textbf{উৎস-সংশোধন (BN-SRC-071):} হিমায়িত ইংরেজি প্রমাণের শেষ ধাপটি নিঃসরণ উপপাদ্যের আরেক প্রয়োগ বলে, কিন্তু $\Gamma \Proves \ltrue \lif \lforall[x][!A(x)]$ থেকে $\Gamma \Proves \lforall[x][!A(x)]$ পেতে $\ltrue$ স্বতঃসিদ্ধটি ও মোডাস পোনেন্স লাগে। বাংলা পাঠে সেই প্রকৃত সমর্থনটি বলা হয়েছে।} % BN-SRC-071 TARGET-CORRECTION-END
\end{proof}

\begin{prop}
\ollabel{prop:provability-quantifiers}
নিচের প্রতিটি বক্তব্যে দেখানো পদটি বদ্ধ:
\begin{tagenumerate}{prvEx,prvAll}
\tagitem{prvEx}{$!A(t) \Proves \lexists[x][!A(x)]$।}{}

\tagitem{prvAll}{$\lforall[x][!A(x)] \Proves !A(t)$।}{}
\end{tagenumerate}
% BN-SRC-072 TARGET-CORRECTION-BEGIN
\par\smallskip\noindent\textnormal{\small\textbf{উৎস-সংশোধন (BN-SRC-072):} হিমায়িত ইংরেজি প্রতিজ্ঞা দুটি $t$-এর কোনো সীমা জানায় না, যদিও যে \olref[qua]{ax:q1} ও~\olref[qua]{ax:q2} থেকে এগুলি পাওয়া হয় সেগুলি কেবল বদ্ধ পদ~$t$-এর জন্য ঘোষিত। বাংলা পাঠে উভয় বক্তব্যের জন্য সেই বদ্ধ-পদ পরিসর স্পষ্ট করা হয়েছে।} % BN-SRC-072 TARGET-CORRECTION-END
\end{prop}

\begin{proof}
\begin{tagenumerate}{prvEx,prvAll}
\tagitem{prvEx}{\olref[qua]{ax:q2} ও নিঃসরণ উপপাদ্য থেকে।}{}
\tagitem{prvAll}{\olref[qua]{ax:q1} ও নিঃসরণ উপপাদ্য থেকে।}{}
\end{tagenumerate}
\end{proof}

\end{document}
